% ===================================================================
%  図表第 16 号 — 大百貨店。— 歳の市。— おもちゃ。
%
%  Ceci n'est PAS une transcription : c'est une traduction, faite en
%  2026, en japonais d'aujourd'hui. Voir texto/ja/10-zuhyo-01.tex
%  pour la consigne, qui vaut pour les seize fichiers.
%
%  Le tableau d'astronomie (bloc 15-1) porte des renvois a LETTRES,
%  (a) a (f), graves sur la planche elle-meme : ils se rangent sous
%  le numero 85, comme le dit gravuri/literi.json. L'ido y saute la
%  parenthese ouvrante du (d) ; la traduction la donne, et le
%  controle passe ce bloc sans rien dire.
%
%  LES TROIS MASQUES GARDENT LEUR NOM, comme le chien Medor et l'ane
%  Aliboron : アルルカン, ポリシネル, ピエロ. Ce sont des personnages,
%  non des objets, et le japonais les nomme ainsi.
% ===================================================================

%%K t16-apar-1 apar
\VUsaut{4.0mm}
\VUtitre{15.0pt}{60}{図表第 16 号}
\VUsaut{2.0mm}
\VUpk{13.2pt}{40}{大百貨店。--- 歳の市。}
\VUpk{13.2pt}{40}{おもちゃ。}
\VUsaut{3.4mm}

%%K t16-01-1 p
1. --- 新しい年が近づいてきた。大きな店はどこも\VUgras{おもちゃ}
\textsuperscript{(1)} と年賀の贈り物の売り場をこしらえた。

%%K t16-01-2 p
私は祖父に、あの立派な陳列を見に歳の市へ連れていってほしいと頼み、祖父は承知して
くれた。

%%K t16-02-1 p
2. --- まず私たちは立派な武器の飾り板の前で足をとめた。\VUgras{小銃}、
\VUgras{軍帽}、\VUgras{サーベル}、\VUgras{剣帯}、\VUgras{肩章}ひとそろいを掛けた
ものもあれば、\VUgras{兜}、\VUgras{胸甲} \textsuperscript{(3)}、\VUgras{ピストル}
\textsuperscript{(4)} を掛けたものもあった。これらはみな\VUgras{幻灯}
\textsuperscript{(2)} よりずっと私の興味を引いた。

%%K t16-tit-1 sub
\VUsaut{3.0mm}
\VUpk{10.2pt}{40}{見ものの数々。}
\VUsaut{2.6mm}

%%K t16-03-1 p
3. --- 同じ売り場に\VUgras{歩哨} \textsuperscript{(5)} が\VUgras{番小屋}に立って
いた。厚紙の動物園をひとつまるごと見張っているように見えた。そこの動物の多いこと。
止まり木の上の\VUgras{おうむ} \textsuperscript{(6)}、鼻に角のある\VUgras{さい}
\textsuperscript{(7)}、\VUgras{となかい} \textsuperscript{(8)}、\VUgras{だちょう}
\textsuperscript{(9)}、くるみをかじる\VUgras{猿} \textsuperscript{(10)}、めんどりを
くわえて運ぶ\VUgras{きつね} \textsuperscript{(11)}、大きなあごを開けた\VUgras{わに}
\textsuperscript{(12)}、\VUgras{らくだ} \textsuperscript{(13)}、姿のよい
\VUgras{グレーハウンド} \textsuperscript{(14)}、\VUgras{ひょう} \textsuperscript{(15)}、
それから私の知らない獣が二匹いて、聞けば\VUgras{いたち} \textsuperscript{(16)} と
\VUgras{ハイエナ} \textsuperscript{(17)} だという。そこには\VUgras{豹}
\textsuperscript{(18)} と\VUgras{おおかみ} \textsuperscript{(21)} もいた。すぐそばには
かわいらしい\VUgras{ねずみ} \textsuperscript{(19)} と、ごく小さなゴムの
\VUgras{はつかねずみ} \textsuperscript{(20)}。最後に\VUgras{かば}
\textsuperscript{(22)} と、こぶのある\VUgras{ひとこぶらくだ} \textsuperscript{(23)} が
この一群をしめくくっていた。そのそばに\VUgras{鉛の兵隊} \textsuperscript{(29)} で
いっぱいの立派な陣営がなかったら、私は何時間でもそこにいたことだろう。それが私の
心を引いてしまったのである。

%%K t16-tit-2 sub
\VUsaut{4.0mm}
\VUpk{10.2pt}{40}{兵隊と戦争。}
\VUsaut{2.8mm}

%%K t16-04-1 p
4. --- 祖父は\VUgras{廃兵} \textsuperscript{(24)} で、\VUgras{傷}のために軍を去ら
ねばならなかった。その傷が\VUgras{武功章}をもたらしたのである。祖父は自分の出た
\VUgras{戦役}の話を私にするのがなによりも好きだ。この小さなミニチュアの軍隊の
さまざまの動きを説明して、たいそう嬉しそうだった。歳の市に来ていた本物の兵隊の一組
——\VUgras{工兵} \textsuperscript{(25)}、\VUgras{ベレー帽}をかぶった\VUgras{アルプス猟兵}
\textsuperscript{(26)}、歩兵の\VUgras{曹長} \textsuperscript{(27)}、袖に金の\VUgras{飾り
モール}をつけた砲兵の\VUgras{軍曹} \textsuperscript{(28)} ——が、その説明を聞こうと
私たちのそばで足をとめた。

%%K t16-05-1 p
5. --- 祖父はこんな立派な聞き手を得てすっかり得意になり、その場で完全な作戦計画を
考え出した。

%%K t16-05-2 p
\VUgras{城壁}と\VUgras{濠}をめぐらした要塞の町は\VUgras{敵軍}に囲まれていた。敵軍は
長いあいだ\VUgras{砲撃}したうえで、いよいよ\VUgras{強襲}にかかろうとしていた。だが
\VUgras{囲まれた者}たちは飢えに追われて、\VUgras{囲む者}たちの\VUgras{陣営}へ
\VUgras{討って出}ようと試みた。\VUgras{天幕}のまわりの頑強な\VUgras{戦闘}でこの
\VUgras{攻撃}を退けると、\VUgras{勝った}囲む者たちは\VUgras{騎兵隊}を放って、
\VUgras{負けた}攻め手を\VUgras{追わ}せた。攻め手は\VUgras{退き}、\VUgras{戦場}を
\VUgras{死者}と\VUgras{負傷者}でおおった。\VUgras{担架兵}がこの後者を
\VUgras{野戦病院}へ運び、\VUgras{外科医}に手当てをさせた。

%%K t16-05-3 p
いくつかの\VUgras{大隊}が\VUgras{方陣}を組んで勇ましく抵抗したにもかかわらず、
\VUgras{敗北}は決定的だった。軍の残りは\VUgras{逃げ}、多くの\VUgras{捕虜}を残した。
敵は町を占領して\VUgras{勝利}をしめくくろうとしていた。おそらくこれが戦役の終わりで
あり、\VUgras{休戦}のあとには\VUgras{講和条約}を結ぶことができるだろう。

%%K t16-tit-3 sub
\VUsaut{4.4mm}
\VUpk{10.2pt}{40}{年賀の贈り物。}
\VUsaut{2.8mm}

%%K t16-06-1 p
6. --- 若い兵隊たちは老兵の絵のような話を笑いながら聞き、さらにいくつか尋ねた。私は
といえば、売り場をまわって立派な\VUgras{弩} \textsuperscript{(32)} と、\VUgras{弓}
\textsuperscript{(33)} とその\VUgras{矢} \textsuperscript{(31)} を眺めた。そこにも
動物の一組があった。二羽の\VUgras{鳴く鳥} \textsuperscript{(34)}、すなわちからくり仕掛けの
\VUgras{ナイチンゲール}と、のどいっぱいに歌う\VUgras{うぐいす}。それに妙な獣の一列
——\VUgras{こうもり} \textsuperscript{(35)}、\VUgras{ボア} \textsuperscript{(36)}、
\VUgras{かめ} \textsuperscript{(37)}、\VUgras{あざらし} \textsuperscript{(38)}、
\VUgras{鯨} \textsuperscript{(39)}、\VUgras{さめ} \textsuperscript{(40)} ——で、みな
金箔の皮でできていた。

%%K t16-07-1 p
7. --- それを長く眺めてはいなかった。かねて夢に見ていたものが目に入ったからである。
みごとな\VUgras{人形芝居の舞台} \textsuperscript{(41)} で、立派な\VUgras{背景}と、
うっとりするような赤い\VUgras{幕}がついている。\VUgras{舞台}の上では二人の役者が、
たいそう面白い\VUgras{芝居}の\VUgras{役}を演じているらしかった。桟敷や\VUgras{肘掛
椅子}に座り、\VUgras{楽長}の後ろの\VUgras{平土間}に陣取った厚紙の小さな
\VUgras{見物人}たちが、しきりに拍手していたからである。

%%K t16-07-2 p
惜しいことに、この芝居小屋はたいそう値が張った。

%%K t16-08-1 p
8. --- そもそもおもちゃを選ぶのはむずかしかった。飾り窓にはありとあらゆるおもちゃが
積み上げてあったからである。\VUgras{アルルカン} \textsuperscript{(42)}、
\VUgras{ポリシネル} \textsuperscript{(43)}、\VUgras{ピエロ} \textsuperscript{(44)}、
翼を打ちふる\VUgras{みみずく} \textsuperscript{(45)}、口笛を吹く\VUgras{くろうたどり}
\textsuperscript{(46)}、吠える\VUgras{プードル}、\VUgras{蓄音機} \textsuperscript{(47)}、
\VUgras{知恵の輪}。そのうちの一つがたいそう面白かった。

%%K t16-noto-1 noto
\VUnotes{143.2mm}{%
\textit{(*) 図表第 6 号を見よ。}%
}

%%K t16-08-1 p suite
それは\VUgras{海戦} \textsuperscript{(48)} を写したもので、\VUgras{ココやし}の茂る陸地の
近くで、一隻の\VUgras{船}が\VUgras{海賊}に襲われているのだった。

%%K t16-09-1 p
9. --- こうした珍しいものをゆっくり眺めることができた。店には人が少なかったからで
ある。ぶらぶら歩きの客がわずかにいるだけで、あとは\VUgras{竜騎兵} \textsuperscript{(49)}
と\VUgras{集金人} \textsuperscript{(50)} がいるくらいだった。この集金人は本
\VUgras{勘定台} \textsuperscript{(51)} で\VUgras{為替手形}を差し出していた。

%%K t16-10-1 p
10. --- 私は祖父に、\VUgras{女会計}の後ろの棚に並べた本は何に使うのかと尋ねた。
祖父は\VUgras{帳簿}だと答えた。

%%K t16-tit-4 sub
\VUsaut{3.6mm}
\VUpk{10.2pt}{40}{店をひとまわり。}
\VUsaut{2.8mm}

%%K t16-11-1 p
11. --- おもちゃの売り場を出て、私たちは馬具の売り場へ\VUgras{鞍} \textsuperscript{(53)}、
\VUgras{あぶみ} \textsuperscript{(54)}、\VUgras{おもがい} \textsuperscript{(55)}、
\VUgras{はみ} \textsuperscript{(56)}、\VUgras{むち}を見に行った。\VUgras{売り場主任}
\textsuperscript{(57)} が祖父にいろいろ説明しているあいだ、私は\VUgras{磁器}と
\VUgras{ガラス器} \textsuperscript{(58)} の陳列を見に行った。そこには立派な
\VUgras{サラダ鉢}と、きゃしゃな\VUgras{茶瓶} \textsuperscript{(59)} があった。

%%K t16-12-1 p
12. --- 二階へ上がるため、私たちは\VUgras{コルセット} \textsuperscript{(60)} の飾り窓の
後ろを通った。靴下売り場のそばで、たいそう立派な\VUgras{ズボン下} \textsuperscript{(63)} が
並べてあった。すぐ近くでは\VUgras{女店員} \textsuperscript{(61)} が\VUgras{女客}
\textsuperscript{(62)} に美しい絹の\VUgras{織物} \textsuperscript{(64)} を選ばせていた。

%%K t16-12-2 p
階段を上がりながら、店が日本の\VUgras{提灯} \textsuperscript{(65)}、\VUgras{日傘}
\textsuperscript{(66)}、それに大きな\VUgras{しゅろ} \textsuperscript{(70)} で飾られている
のに気づいた。

%%K t16-13-1 p
13. --- 二階には見るものがあまりなかった。家具(*)、\VUgras{金庫} \textsuperscript{(67)}、
\VUgras{たんす} \textsuperscript{(68)}、\VUgras{乳母車} \textsuperscript{(69)} だけである。
けれども店全体を見わたす眺めはたいそう\VUgras{面白かった}。

%%K t16-14-1 p
14. --- 足もとには楽器が見えた。\VUgras{たてごと} \textsuperscript{(71)}、
\VUgras{ギター} \textsuperscript{(72)}、\VUgras{マンドリン} \textsuperscript{(73)}、
\VUgras{ピストン付きコルネット} \textsuperscript{(74)}、\VUgras{クラリネット}
\textsuperscript{(75)}、\VUgras{弓} \textsuperscript{(76)} を添えたバイオリン、それに
\VUgras{大太鼓} \textsuperscript{(77)} まであった。もう少し先の紙の売り場では、
\VUgras{学生} \textsuperscript{(78)} が\VUgras{店員} \textsuperscript{(79)} を相手に紙挟みの
値を掛け合っていた。そのそばに立派な\VUgras{いろは読本} \textsuperscript{(80)} が見えた。

%%K t16-14-2 p
店の真ん中には高い\VUgras{クリスマスの木} \textsuperscript{(81)} が立てられ、ろうそくや
小間物やありとあらゆるおもちゃ——\VUgras{木馬} \textsuperscript{(82)}、\VUgras{人形}
\textsuperscript{(83)} など——がいっぱいに掛かっていた。

%%K t16-14-3 p
\VUgras{香水売り場} \textsuperscript{(84)} は\VUgras{歯みがき}やいろいろの\VUgras{香水}
を置いていて、たいそうよい香りを放ち、それが私たちのところまで漂ってきた。

%%K t16-tit-5 sub
\VUsaut{3.4mm}
\VUpk{10.2pt}{40}{私たちは買い物をする。}
\VUsaut{2.8mm}

%%K t16-15-1 p
15. --- 祖父が長すぎると思いはじめたので、私たちは大階段を下りた。祖父は
\VUgras{天文の掛け図} \textsuperscript{(85)} の前でしばらく足をとめた。そこに描いてあるのは
——と祖父は言った——\VUgras{彗星} \textsuperscript{(a)}、\VUgras{月食と日食}
\textsuperscript{(b)}、\VUgras{四つの方位} \textsuperscript{(c)} \VUgras{（北、南、東、西）}
をしるした風配図、二つの\VUgras{半球} \textsuperscript{(d)} の図、\VUgras{惑星}
\textsuperscript{(e)}、\VUgras{太陽系} \textsuperscript{(f)} である。私にはさっぱり分からな
かった。それより、通りかかった\VUgras{配達の若い者} \textsuperscript{(86)} の飾りモールの
ついたお仕着せや、そばにあった美しい\VUgras{扇} \textsuperscript{(87)}、\VUgras{がま口}
\textsuperscript{(88)}、\VUgras{紙入れ} \textsuperscript{(89)} のほうがずっと面白かった。

%%K t16-16-1 p
16. --- 陳列台の上の\VUgras{羽根飾り} \textsuperscript{(90)} や\VUgras{帯留めの留め金}
\textsuperscript{(91)}、それに籠に品よく生けた美しい\VUgras{造花} \textsuperscript{(94)} にも
見とれた。\VUgras{すみれ}、\VUgras{においあらせいとう}、\VUgras{ヒヤシンス}
\textsuperscript{(95)}、\VUgras{ゼラニウム} \textsuperscript{(96)}、\VUgras{ゆり}
\textsuperscript{(97)}、\VUgras{きんれんか} \textsuperscript{(98)} は、たいそうよくできて
いた。

%%K t16-tit-6 sub
\VUsaut{4.4mm}
\VUpk{10.2pt}{40}{祖父の贈り物。}
\VUsaut{2.8mm}

%%K t16-17-1 p
17. --- 私は祖父に、\VUgras{歯ブラシ} \textsuperscript{(92)} を一本買ってほしいと頼んだ。
\VUgras{ヘアピン} \textsuperscript{(93)} のそばの箱に入っていたものである。祖父は承知して
くれ、私は自分で勘定台へ払いに行った。そのそばでは大きな\VUgras{送り荷}
\textsuperscript{(100)} がこしらえられていた。ある\VUgras{客} \textsuperscript{(99)} あての
ものである。

%%K t16-18-1 p
18. --- 突然、美術の\VUgras{ブロンズ} \textsuperscript{(101)} のそばに立っていた人たちの
あいだに、ちょっとした騒ぎが起こった。\VUgras{泥棒} \textsuperscript{(102)} が誰かの物を
盗もうとしているところを、\VUgras{監視係} \textsuperscript{(103)} に取り押さえられたので
ある。人をやって巡査を呼び、\VUgras{捕らえ}させた。私たちが店を出るちょうどそのとき、
その罪人は留置場へ引かれていくところだった。
\VUsaut{6.0mm}
\VUfilet{22mm}
